\section{Open questions}\label{sec:oq}

The following five questions remain genuinely open after this replication. Each
targets a specific unverified aspect of the paper and has a concrete,
free-endpoint next-step probe.

\begin{enumerate}
\item \textbf{Do the derivative recipes hold at second order (Hessian, $\alpha_{ZZ}$)?}
  \emph{Basis:} Only first-order derivatives were exercised. The paper's C5
  (Newton's method) and Fig.\ 5 polarizability curves depend on the Hessian
  recipe. The observed $10^3\times$ HF-vs-FD accuracy penalty at first order
  raises concern about second-order compounding.
  \emph{Next steps:} Add a second finite-difference layer to compute
  $d^2 E / d R^2$ via VQE-FD and via the paper's Eq.\ 6--7 Hessian formula;
  compare against PySCF FCI Hessian. Implement $\alpha_{ZZ} = d^2 E / dF_Z^2$
  under a static field perturbation and compare to the paper's Fig.\ 5.

\item \textbf{Does the framework survive adaptive ansatze (ADAPT-VQE)?}
  \emph{Basis:} The paper uses a fixed UCCSD-restricted ansatz. Post-2022 the
  practical VQE default is ADAPT-VQE (adds/removes operators during
  optimization) which raises parameter-shift-continuity questions the paper
  does not address.
  \emph{Next steps:} Rerun H$_2$/STO-3G with PennyLane's AdaptiveOptimizer +
  \texttt{qml.qchem.excitations} operator pool; at each ADAPT growth step
  compute VQE-FD, VQE-HF, and PySCF-FCI gradients and check consistency.

\item \textbf{At what shot budget does VQE-FD lose its advantage over Hellmann--Feynman?}
  \emph{Basis:} The $10^3\times$ FD-vs-HF gap here is a statevector-simulator
  artifact. Under finite shots, VQE-FD's doubled variance (two shifted
  expectations subtracted) will dominate. The paper's \S III.C stops short of
  a shot-noise crossover analysis.
  \emph{Next steps:} Rerun the 5-point scan with shot counts in
  $\{10^3, 10^4, 10^5, 10^6\}$; measure gradient variance for both recipes and
  locate the crossover. Optionally add a depolarizing-noise model at
  $p \in \{10^{-4}, 10^{-3}\}$.

\item \textbf{Does the paper's TS-search recipe reproduce the H+H$_2$ symmetric saddle at (0.936, 0.936) \AA{} via P-RFO on the VQE Hessian?}
  \emph{Basis:} Saddle-point search (paper C7) was not exercised here. It
  requires a reliable Hessian, eigenvector-following, and initial-guess
  sensitivity that gradient-only opt does not face.
  \emph{Next steps:} Extend to a 3-atom H$_3$/STO-3G Hamiltonian
  ($\sim 6$--8 qubits) via \texttt{qml.qchem.molecular\_hamiltonian}; write a
  P-RFO driver on the existing gradient+Hessian VQE stack; seed from
  $(0.936, 0.936)$ \AA{} and verify convergence to a first-order saddle with
  one imaginary mode along the reaction coordinate.

\item \textbf{Does the recipe compose with a QM/MM or DMET embedding?}
  \emph{Basis:} The paper treats isolated H$_2$ only. In embedded settings the
  environment-derived integrals are also geometry-dependent, introducing a
  Pulay-force-analog term the paper's Hellmann--Feynman formula does not
  account for.
  \emph{Next steps:} Build a minimal H$_2$-in-point-charge-field toy embedding
  (two static $\pm 0.1e$ charges at fixed positions), compute
  $dE_{\text{VQE}}/dR_{\rm H}$ with the field-modified Hamiltonian, compare
  against numerical gradient of the same setup, and quantify the systematic
  shift from ignoring $dH_{\rm environment}/dR$.
\end{enumerate}
